\section{Preliminaries}\label{sect2}
The results in the section for which proofs are omitted are known;
proofs for them are given 
in~\cite{AFM}, \cite{BZ}, or~\cite{Blau2}. Throughout,
$(A, B)$ is a table algebra (TA).



There is a positive definite Hermitian form $(\ ,\ )$ on $A$ such that
for
all $b,c,d \in B$,
\begin{center}
$(b,c) = \beta_{bc^*1}$; and
$(bc,d) = (b,dc^*) = (c,b^*d)$.
\end{center}

\begin{defi}
\label{def:TA-hom}
Let $(A, B), (U, V)$ be TAs. A \emph{table algebra} 
\textup{(}TA-\textup{)}\emph{homomorphism} 
$\psi: (A, B) \rightarrow (U, V)$ 
is an algebra homomorphism $\psi: A \rightarrow U, (\psi (1_A) = 1_U)$,
such that for each $b \in B$, there is some $v \in V$ and $\alpha_{b}
\in \R_{>0}$ with
$ \psi(b) = \alpha_{b}v$.
Define
$V_{\psi} \coloneqq \{ v \in V \mid \psi(b) = \alpha_bv \mbox{ for some }b \in
B \mbox{ and } \alpha_b \in \R_{>0}\}$. Then $V_{\psi}$ is a closed
subset of $V$. 
Now $\psi$ is called a TA-\emph{isomorphism} if $\psi$ is one-to-one
and $V_{\psi} = V$. This is equivalent to $\psi: A \rightarrow U$
being an algebra isomorphism. If there exists such a TA-isomorphism, 
we 
write $(A, B) \cong (U, V)$.
\end{defi}

\begin{defi}
\label{def:kernel}
The \emph{kernel} of a TA-homomorphism $\psi$ is defined as
$\ker \psi \coloneqq \{b \in B \vert \psi(b) = \alpha_b1_U \mbox{ for some
}\alpha_b \in \R_{>0}\}$.
\end{defi}

Then $\psi$ preserves the respective anti-automorphisms, $\ker \psi$ is a normal closed subset of $B$, and the
following ``Fundamental Homomorphism Theorem'' holds:
 

\begin{prop}
\label{prop:FHT}
Let $\psi : (A, B) \rightarrow (U, V)$ be a TA-epimorphism of STAs. Let
$C = \ker \psi$. Let $e = e_C \coloneqq o(C)^{-1}C^+$, a central
idempotent of $A$. Then $\pi : (A, B) \rightarrow (A//C, B//C)$, where
$\pi(b) = be$ for all $b \in B$, is a TA-epimorphism, as $be =
(\delta(b)/\delta(b//C))b//C$. Furthermore, there is a TA-isomorphism
$\overline{\psi} : (A//C, B//C) \rightarrow (U, V)$ such that
$\overline{\psi}\circ \pi = \psi$. Finally, since both $(A//C, B//C)$
and $(U, V)$ are standard, then for all $b \in B$,
$\overline{\psi}(b//C) = v$, where $\psi(b) = \alpha_bv$ for some
$\alpha_b \in {\R}_{>0}$. 
\end{prop} 


Proposition~\ref{prop:FHT} has the following consequence:

\begin{lemm}
\label{lem:isom}
Let $(A, B)$ be a STA with closed subsets $C, D$ where $D$ is normal in
$B$. Then $DC = CD$ is a closed subset, $D\cap C$ is a normal closed
subset of $C$, and
\[
C//(D\cap C) \cong CD//D,
\] 
via a TA-isomorphism that yields the correspondence $c//D\cap C
\leftrightarrow c//D$, for all $c \in C$.
\end{lemm}


The first two parts of the next proposition are contained in~\cite
[Proposition 2.13]{BZ}, and part~\ref{prop2.5_iii} is proved for hypergroups in~\cite
[Lemma 4.6]{FZ}. Parts~\ref{prop2.5_iii}, \ref{prop2.5_iv}, and~\ref{prop2.5_v} are proved below.

\begin{prop}
\label{prop:subset_corresp}
Let $(A, B)$ be a STA and $C$ a closed subset of $B$. Then the
following hold:
\begin{enumerate}[label = \textup{(}\roman*\textup{)}] 
\item\label{prop2.5_i} The correspondence $D \mapsto D//C$ is a bijection between the set of
closed subsets of $B$ that contain $C$ and the set of closed subsets of
$B//C$.
\item\label{prop2.5_ii} Suppose that $D$ is a closed subset of $B$ with $C \subseteq D \subseteq
B$. Then for all $b \in B$, $(b//C)//(D//C) = b//D$. Hence,
$(B//C)//(D//C) = B//D$.
\item\label{prop2.5_iii} Suppose that $C \subseteq D \subseteq B$. Then $D$ is strongly normal
in $B$ if and only if $D//C$ is strongly normal in $B//C$.
\item\label{prop2.5_iv} Suppose that $C \subseteq D \subseteq B$ and $D$ is normal in $B$. Then
$D//C$ is normal in $B//C$.
\item\label{prop2.5_v} Suppose that $C \subseteq D \subseteq B$ and $C$ is normal in $B$. Then
$D$ is a normal subset of $B$ if and only if $D//C$ is a normal subset
of $B//C$.
\end{enumerate}
\end{prop}
